% GENERATED FILE. Edit canonical lessons, then run npm run generate:books.
% canonical-source-hash: fnv1a64:461bc5fc67878093
% canonical-lessons: ML-C93-kal, ML-C93-angal, ML-C93-counted, ML-R93-plural-recall

\chapter{More Than One}
\label{ch:ml-plural}
\hlchaptermodality{93}

\begin{chapteropening}
\textbf{By the end of this chapter:} \emph{I can say that there is more than one of something, and I know when to leave the ending off.}

From cold: the plural ending on a plain noun, the shape it takes on a noun ending in the anusvara, and the case where a number in front makes the ending unnecessary.
\end{chapteropening}

\section[kuṭṭikaḷ]{\ml{കുട്ടികൾ} (kuṭṭikaḷ) — children}

\label{lesson:ML-C93-kal}



\begin{quote}
Say \emph{a child}. Say \emph{a shop}. Both are one of the thing.
\end{quote}

\begin{grammarlens}[title={Grammar Lens: the ending goes on the end}]
Every noun you own says \emph{one}. To say \emph{more than one}, put \textbf{\ml{കൾ}} on the end:

\noindent
\begin{tabularx}{\linewidth}{@{}>{\raggedright\arraybackslash}X>{\raggedright\arraybackslash}X@{}}
\toprule
\textbf{one} & \textbf{more than one} \\
\midrule
\textbf{\ml{കുട്ടി}} \emph{kuṭṭi} — a child & \textbf{\ml{കുട്ടികൾ}} \emph{kuṭṭikaḷ} — children \\
\textbf{\ml{കട}} \emph{kaṭa} — a shop & \textbf{\ml{കടകൾ}} \emph{kaṭakaḷ} — shops \\
\bottomrule
\end{tabularx}

\textbf{The front of the word does not move.} Nothing is inserted, nothing is replaced; the ending arrives on the end and the word you already had is still sitting there in front of it.

That is the same shape as everything Malayalam has asked of you so far. The \emph{in} ending, the \emph{to} ending, the \emph{of} ending, the \emph{if} ending — each one came onto the back of a word and left the front of it alone. \textbf{A Malayalam word grows rightward}, and this is one more thing growing there.
\end{grammarlens}

\begin{grammarlens}[title={What you've built}]
You can now point at more than one of anything you can name, which is most of what this book has taught you. \textbf{Two nouns became four in one line}, and the same line will do it to every other noun you own.

English gets this wrong often enough to have a page of exceptions for it — \emph{child} and \emph{children} share four letters and part company after them. Here the word you learned stays whole and visible.
\end{grammarlens}

\subsection*{Your turn}
\begin{itemize}
\raggedright
  \item \emph{Say it:} \emph{kuṭṭi}, then \emph{kuṭṭikaḷ}
  \item \emph{Say it:} \emph{kaṭa}, then \emph{kaṭakaḷ}
  \item \emph{Look:} at \textbf{\ml{കുട്ടികൾ}} and find where \textbf{\ml{കുട്ടി}} ends
  \item \emph{Say it:} which end of the word the new part arrived on
\end{itemize}

\begin{tcolorbox}[breakable,colback=teal!4,colframe=teal!35!black,title={Before you move on}]
What makes a noun more than one? (\textbf{\ml{കൾ} on the end}.) Does anything change at the front of the word? (\textbf{No.})
\end{tcolorbox}
\section[pustakaṅṅaḷ]{\ml{പുസ്തകങ്ങൾ} (pustakaṅṅaḷ) — books}

\label{lesson:ML-C93-angal}



\begin{quote}
Say \emph{a book}. Say what its last letter is. Then make \emph{a child} into \emph{children}.
\end{quote}

\begin{grammarlens}[title={Grammar Lens: a noun ending in \ml{ം} trades it away}]
\textbf{\ml{പുസ്തകം}} does not become \emph{pustakaṁkaḷ}. The \textbf{\ml{ം}} comes off and \textbf{\ml{ങ്ങൾ}} goes on in its place:

\noindent
\begin{tabularx}{\linewidth}{@{}>{\raggedright\arraybackslash}X>{\raggedright\arraybackslash}X@{}}
\toprule
\textbf{one} & \textbf{more than one} \\
\midrule
\textbf{\ml{പുസ്തകം}} \emph{pustakaṁ} — a book & \textbf{\ml{പുസ്തകങ്ങൾ}} \emph{pustakaṅṅaḷ} — books \\
\textbf{\ml{മരം}} \emph{maraṁ} — a tree & \textbf{\ml{മരങ്ങൾ}} \emph{maraṅṅaḷ} — trees \\
\bottomrule
\end{tabularx}

The money lesson said that a word ending in \textbf{\ml{ം}} tells you in advance how it will behave. \textbf{This is one of the things it was telling you.} You do not have to remember which nouns take which ending — you can see it on the page, in the last letter.
\end{grammarlens}

\begin{grammarlens}[title={What you've built}]
Now look at a word you have had since your first pages.

\begin{quote}
\textbf{\ml{നിങ്ങൾ}} \emph{niṅṅaḷ} — you.
\end{quote}

\textbf{Read the end of it.} That is the ending you learned a moment ago, sitting where it has always sat. The polite \emph{you} was introduced to you as the plural \emph{you} used for one person — and this is what made it plural.

Nothing you were taught then was wrong, and nothing has to be unlearned. \textbf{A piece of the language you had been pronouncing for months turns out to have a name and a job}, and you can now put it on nouns of your own.
\end{grammarlens}

\subsection*{Your turn}
\begin{itemize}
\raggedright
  \item \emph{Say it:} \emph{pustakaṁ}, then \emph{pustakaṅṅaḷ}
  \item \emph{Say it:} \emph{maraṁ}, then \emph{maraṅṅaḷ}
  \item \emph{Say it:} \emph{kuṭṭi}, then \emph{kuṭṭikaḷ}
  \item \emph{Look:} at \textbf{\ml{നിങ്ങൾ}} and \textbf{\ml{പുസ്തകങ്ങൾ}}, and find the part they share
  \item \emph{Say it:} which letter on the end of a noun sends you to this ending
\end{itemize}

\begin{tcolorbox}[breakable,colback=teal!4,colframe=teal!35!black,title={Before you move on}]
What happens to the \textbf{\ml{ം}} when the plural arrives? (\textbf{It is traded for \ml{ങ്ങൾ}}.) Which noun tells you to expect that? (\textbf{One whose last letter is \ml{ം}}.) Where else have you been saying that ending? (\textbf{On the end of \emph{you} and \emph{we}}.)
\end{tcolorbox}
\section[raṇṭŭ pustakaṁ]{\ml{രണ്ട്} \ml{പുസ്തകം} (raṇṭŭ pustakaṁ) — two books}

\label{lesson:ML-C93-counted}



\begin{quote}
Make \emph{a book} into \emph{books}. Then say \emph{fifty rupees}.
\end{quote}

\begin{grammarlens}[title={Grammar Lens: the number has said it already}]
\begin{quote}
\textbf{\ml{രണ്ട്} \ml{പുസ്തകം}} — \emph{raṇṭŭ pustakaṁ} — \textbf{two books.}
\end{quote}

The noun is back to the shape it had before this chapter started, and the meaning is still more than one. \textbf{\ml{രണ്ട്}} in front has done the counting, so the ending has nothing left to add and is left off.

\noindent
\begin{tabularx}{\linewidth}{@{}>{\raggedright\arraybackslash}X>{\raggedright\arraybackslash}X@{}}
\toprule
\textbf{} & \textbf{} \\
\midrule
\textbf{\ml{പുസ്തകങ്ങൾ}} & books — how many is not said \\
\textbf{\ml{രണ്ട്} \ml{പുസ്തകം}} & two books — the number says how many \\
\bottomrule
\end{tabularx}

Both are good Malayalam. The ending is carrying information, and where the number already carries it, the ending can go.
\end{grammarlens}

\begin{grammarlens}[title={What you've built}]
This is not new to you. Go back over your own sentences:

\begin{quote}
\textbf{\ml{അമ്പത്} \ml{രൂപ}} — fifty rupees.
\end{quote}

\begin{quote}
\textbf{\ml{നൂറ്} \ml{രൂപ}} — a hundred rupees.
\end{quote}

\begin{quote}
\textbf{\ml{രണ്ട്} \ml{മണി}} — two o'clock.
\end{quote}

\textbf{Not one of those nouns took an ending}, and every one of them was right. You had been counting this way for chapters before anybody handed you the ending, because the shape was correct and you had no reason to question it.

So the lesson here costs you nothing to do and something to notice: \textbf{an English speaker's instinct to reach for the plural after a number is the habit to put down.} Say the number, then say the noun as you learned it.
\end{grammarlens}

\subsection*{Your turn}
Say these aloud:

\begin{itemize}
\raggedright
  \item \emph{pustakaṅṅaḷ}, then \emph{raṇṭŭ pustakaṁ}
  \item \emph{ampathŭ rūpa}
  \item which of the two ways says how many
  \item what you would put after \emph{mūnnŭ} to mean three books
\end{itemize}

\begin{tcolorbox}[breakable,colback=teal!4,colframe=teal!35!black,title={Before you move on}]
Does a noun with a number in front need the plural ending? (\textbf{No} — the number has said how many.) Had you been doing that already? (\textbf{Yes}, every time you gave a price or an hour.)
\end{tcolorbox}
\section[bahuvacanam enna ōrmma]{(more than one) — from cold}

\label{lesson:ML-R93-plural-recall}



\begin{quote}
Close the lessons before this one. Make \emph{a child} into \emph{children}, and \emph{a book} into \emph{books}. The two answers do not end the same way.
\end{quote}

\begin{grammarlens}[title={Grammar Lens: one ending, wearing two faces}]
\noindent
\begin{tabularx}{\linewidth}{@{}>{\raggedright\arraybackslash}X>{\raggedright\arraybackslash}X@{}}
\toprule
\textbf{one} & \textbf{more than one} \\
\midrule
\textbf{\ml{കുട്ടി}} \emph{kuṭṭi} & \textbf{\ml{കുട്ടികൾ}} \emph{kuṭṭikaḷ} \\
\textbf{\ml{കട}} \emph{kaṭa} & \textbf{\ml{കടകൾ}} \emph{kaṭakaḷ} \\
\textbf{\ml{പുസ്തകം}} \emph{pustakaṁ} & \textbf{\ml{പുസ്തകങ്ങൾ}} \emph{pustakaṅṅaḷ} \\
\textbf{\ml{മരം}} \emph{maraṁ} & \textbf{\ml{മരങ്ങൾ}} \emph{maraṅṅaḷ} \\
\bottomrule
\end{tabularx}

\textbf{The last letter of the noun decides which face you see.} A noun ending in \textbf{\ml{ം}} trades that letter for \textbf{\ml{ങ്ങൾ}}; anything else takes \textbf{\ml{കൾ}} on the end and keeps everything it had.

Look down the right-hand column and read only the ends: \textbf{\ml{ൾ}} closes all four. The two shapes are one ending, and the \textbf{\ml{ം}} is what bends it.
\end{grammarlens}

\begin{grammarlens}[title={What you've built}]
Two things worth holding onto, and neither is a new word.

\textbf{The ending was already in your mouth.} \textbf{\ml{നിങ്ങൾ}} — the respectful \emph{you} — has carried it since your earliest pages, and it was plural in shape from the day you met it.

\textbf{The ending comes off again when a number is there.}

\begin{quote}
\textbf{\ml{രണ്ട്} \ml{പുസ്തകം}} — two books. \textbf{\ml{അമ്പത്} \ml{രൂപ}} — fifty rupees.
\end{quote}

The number in front has said how many, so the noun goes back to the form you first learned. Every price you have given and every hour you have told was already built this way.

So this chapter added \textbf{no vocabulary at all}. What it added was a reading of material you already had: an ending you had been pronouncing, and a habit you had been performing correctly without being told why.
\end{grammarlens}

\subsection*{Your turn}
\begin{itemize}
\raggedright
  \item \emph{Say it:} \emph{kuṭṭi} $\to$ \emph{kuṭṭikaḷ}
  \item \emph{Say it:} \emph{pustakaṁ} $\to$ \emph{pustakaṅṅaḷ}
  \item \emph{Say it:} \emph{maraṁ} $\to$ \emph{maraṅṅaḷ}
  \item \emph{Say it:} \emph{raṇṭŭ pustakaṁ}
  \item \emph{Look:} at \textbf{\ml{നിങ്ങൾ}}, and say what its last part is doing
\end{itemize}

\begin{tcolorbox}[breakable,colback=teal!4,colframe=teal!35!black,title={Before you move on}]
Which ending does a noun ending in \textbf{\ml{ം}} take? (\textbf{\ml{ങ്ങൾ}}, in place of the \textbf{\ml{ം}}.) Which does every other noun take? (\textbf{\ml{കൾ}}, on the end.) And when a number stands in front? (\textbf{Neither} — the noun stays as it was.)
\end{tcolorbox}
